%%%%%%%%%%%%%%%%%%%%%%%%%%%%%%%%%%%%%%%%%%%%%%%%%%%%%%%%%%%%%%%%%%%%%%%%%%%%%%
%% reference:
%% http://www-i6.informatik.rwth-aachen.de/~dreuw/latexbeamerposter.php
%%%%%%%%%%%%%%%%%%%%%%%%%%%%%%%%%%%%%%%%%%%%%%%%%%%%%%%%%%%%%%%%%%%%%%%%%%%%%%

\documentclass[final]{beamer}
\mode<presentation>
{
  %\usetheme{I6dv}
  \usetheme{l6dv-2}
  %\usetheme{l6dv-3}
%  \usetheme{Aachen}
%  \usetheme{Oldi6}
%  \usetheme{I6td}
  %\usetheme{I6dv}
%  \usetheme{I6pd}
%  \usetheme{I6pd2}
}
% additional settings
\usefonttheme[onlymath]{serif}
\setbeamerfont{itemize}{size=\normalsize}
\setbeamerfont{itemize/enumerate body}{size=\normalsize}
\setbeamerfont{itemize/enumerate subbody}{size=\small}
\setbeamertemplate{bibliography entry title}{}
\setbeamertemplate{bibliography entry location}{}
\setbeamertemplate{bibliography entry note}{}

\usepackage{natbib}

%\renewcommand\printleftlogo
%  {\begin{center}
%     \resizebox{\textwidth}{!}%
%       {\includegraphics{GMU_Logo.eps}}
%   \end{center}
%  }

% additional packages
\usepackage{times}
\usepackage{amsmath,amsthm, amssymb, latexsym,epsfig,graphicx}
\usepackage{mathtools}
\usepackage{exscale}
\usepackage{multirow}
%\usepackage{subfig}
\usepackage{color}

\usepackage{pgfplots}
\pgfplotsset{every axis/.append style={
                    xlabel={$x$},          % default put x on x-axis
                    ylabel={$y$},          % default put y on y-axis
                    label style={font=\tiny},
                    tick label style={font=\tiny}
            }}
\pgfplotsset{every tick label/.append style={font=\tiny}}
\pgfplotsset{every axis/.append style={font=\tiny}}
\pgfplotsset{compat=1.12}



\usepackage{subcaption}
\captionsetup{compatibility=false}
\usepackage{booktabs, array}
%\usepackage{tabularx}
%\usepackage{rotating} %sideways environment
\usepackage[english]{babel}
%\usepackage[latin1]{inputenc}
%\usepackage[orientation=landscape,size=custom,width=200,height=120,scale=1.9]{beamerposter}
\usepackage[size=custom, width=48,height=36,scale=1.4]{beamerposter}
\listfiles


\title{\huge Root Finding with Chebyshev Polynomials in Two Dimensions}
\author[Bouck, Bell]{Lucas Bouck and Ian Bell}
\institute[George Mason University]{George Mason University, Fairfax, VA, U.S.A. and National Institute of Standards and Technology, Boulder, CO, U.S.A}
%\date[Aug. 31 , 2007]{Aug. 31 , 2007}

% abbreviations
\usepackage{xspace}
\makeatletter
%\DeclareRobustCommand\onedot{\futurelet\@let@token\@onedot}
%\def\@onedot{\ifx\@let@token.\else.\null\fi\xspace}
%\def\eg{{e.g}\onedot} \def\Eg{{E.g}\onedot}
%\def\ie{{i.e}\onedot} \def\Ie{{I.e}\onedot}
%\def\cf{{c.f}\onedot} \def\Cf{{C.f}\onedot}
%\def\etc{{etc}\onedot}
%\def\vs{{vs}\onedot}
%\def\wrt{w.r.t\onedot}
%\def\dof{d.o.f\onedot}
%\def\etal{{et al}\onedot}
\makeatother

%%%%%%%%%%%%%%%%%%%%%%%
% Logo for Poster
%%%%%%%%%%%%%%%%%%%%%
%\printleftlogo
%\leftlogo[1]{GMU_Logo.eps}  % defines logo to left of title (with scale factor)
%\rightlogo[0.7]{nsf_logo.jpg}  % same but on right

%%%%%%%%%%%%%%%%%%%%%%%%%%%%%%%%%%%%%%%%%%%%%%%%%%%%%%%%%%%%%%%%%%%%%%%%%%%%%%%%%%%%%%%%%%%%%%%%%%%%%%%%%%%%
%%%%%%%%%%%%%%%%%%%%%%%%%%%%%%%%%%%%%%%%%%%%%%%%%%%%%%%%%%%%%%%%%%%%%%%%%%%%%%%%%%%%%%%%%%%%%%%%%%%%%%%%%%%%
\begin{document}
%\addtobeamertemplate{headline}{} 
%{
%\begin{tikzpicture}[remember picture,overlay]
%\node [shift={(-5.5 cm,-4cm)}] at (current page.north east) {\includegraphics[height=6.5cm]{GMU_Logo}}; 
%\end{tikzpicture} 
%}
\vspace*{-10mm}
\begin{frame}{}
  \begin{columns}[t]
    \hspace*{5mm}

    \begin{column}{.31\linewidth}

    %%%%%%%%%%%%%%%%%%%%%%%%%%%%%%%%%%%%%%%%%%%%%%%%%%%%%%%%%%%%%%%%%%%%%%%%%%%%%%%%%%%%%%%%%%%%%%%%%%%%%%%%%%%%
\begin{block}{Summary}
We want to find all isolated, common roots of $f,g: \mathbb{R}^2\to\mathbb{R}$ in a bounded domain. We replicate the root finding algorithm found in \cite{nakatsukasa_2013}, which approximates $f,g$ with Chebyshev polynomials and finds the roots of the polynomials. We are working to extend the algorithm to non-rectangular domains.
 \end{block}
 \vspace*{2mm}
 \begin{block}{Background, Applications, and Motivation}
 {\color{numhypRed} Background and Applications}
\begin{itemize}
\item {\tt ChebTools} \cite{bell_2018} is a C++ library developed by Ian Bell
\item Inspired by the Matlab library {\tt chebfun} \cite{driscoll_2014}
\item Determine thermodynamic properties of a fluid \cite{kunick_2008}
\end{itemize}
{\color{numhypRed} Motivating Example}\\
$f(x,y)=x^2-y^2, g(x,y)=\sin\left(\pi\sqrt{x^2+y^2}\right)$ with the domain $[-1,1]^2$
\begin{figure}[H] 
\begin{center}
\begin{minipage}{.55\textwidth}
\includegraphics[trim=1cm 0cm 0cm 0cm,width=\textwidth]{surface_plots1.pdf}
\end{minipage}%
\begin{minipage}{.45\textwidth}
\includegraphics[width=\textwidth]{newton_basins.pdf}
\end{minipage}%
\end{center}
\caption{Newton's method would require multiple runs to find all roots, which may not be robust for our application.}
\end{figure}
\end{block}

\begin{block}{Chebyshev Polynomials}
\begin{minipage}{.578\linewidth}
{\color{numhypRed} Definition}
$$T_n(x)=\cos\left(n\arccos(x)\right),\, x\in[-1,1]$$
{\color{numhypRed} Interpolation Convergence \cite[Ch.\ 7, 8]{trefethen_2012}}
\begin{itemize}
\item $f^{(m)}$ has bounded variation $\implies$ polynomial convergence of order $m$
\item $f$ is analytic $\implies$ geometric convergence 
\end{itemize}
\end{minipage}%}%
\hfill%
%\fbox{%
\begin{minipage}{.382\linewidth}
\begin{center}
\begin{figure}
\includegraphics[width=\textwidth]{cheb_poly_plot.pdf}
\end{figure}
\end{center}
\end{minipage}%}
\end{block}
\begin{block}{Summary of Algorithm from \cite{nakatsukasa_2013}}
\begin{itemize}
\item Construct interpolations, $p_f,p_g$, of $f,g$ \cite[Sec.\ 2]{townsend_2013}
\item Apply the B\'{e}zout Resultant Method for $x$ values
\item Apply 1-D root finding techniques for $y$ values
\end{itemize}
{\color{numhypRed} B\'{e}zout Resultant Method}\\
With $p_f,p_g$ of degree $(m_f,n_f),(m_g,n_g)$, we can construct a square {\bf B\'{e}zout Matrix Polynomial} of degree $M\leq m_f+m_g$ and size $n = \max(n_f,n_g)$:
$$B(x) = \sum_{i=0}^{M}B_iT_i(x),$$
\begin{itemize}
\item $\det\left(B(x_0)\right)=0\iff p_f(x_0,\cdot)$ and $p_g(x_0,\cdot)$ have a common root
\item Solving $\det\left(B(x_0)\right)=0$ involves linearizing $B(x)$ 
\end{itemize}
\end{block}



\end{column}
%%% end of first column %%%%%%%%%%%%

%% 2/19/2012 start next column next time %%%%%%%%%%%
%%%%%%%%%%%%%%%%%%%%%%%%%%%%%%%%%%%%%%%%%%%%%%%%%%%%%%%%%%%%%%%%%%%%%%%%%%%%%%%%%%%%%%%%%%%%%%%%%%%%%%%%%%%
%%% begining of second column %%%%%%%%%%%%%%
    \begin{column}{.31\linewidth}
    %%%%%%%%%%%%%%%%%%%%%%%%%%%%%%%%%%%%%%%%%%%%%%%%%%%%%%%%%%%%%%%%%%%%%%%%%%%%%%%%%%%%%%%%%%%%%%%%%%%%%%%%%%%%%%%%%%%%%%%%%%%%%%%%%%%%%%%%%%%

%

\begin{block}{Challenges with the Algorithm}
{\color{numhypRed}Poor Conditioning}\\
\begin{center}
\begin{figure}
\includegraphics[width=.6835\textwidth]{bezout_det_plot.pdf}
\caption{$\det(B(x))$ from our example. The B\'{e}zout technique can square the condition number. \cite[Sec.\ 5]{nakatsukasa_2013}
}
\end{figure}
\end{center}
{\color{numhypRed}Poor Scaling}\\
\begin{itemize}
\item Have to solve a generalized eigenvalue problem \cite[Sec.\ 3]{nakatsukasa_2013}
\item Computational complexity of $\mathcal{O}(M^3n^3)$ \cite[Sec.\ 4]{nakatsukasa_2013}
\end{itemize}
\end{block}

\begin{block}{Our Implementation}
{\color{numhypRed} Refinement}\\
\begin{itemize}
\item {\tt chebfun2} refines roots by repeating the B\'{e}zout method in a smaller region surrounding the root \cite[Sec.\ 7]{nakatsukasa_2013}
\item We refine the roots with Newton's method
\end{itemize}
{\color{numhypRed} 1-D Root Finding}\\
\begin{itemize}
\item Our 1-D algorithm is a mix of bisection and secant methods
\item Faster than companion matrix methods \cite{boyd_2013}
\end{itemize}
{\color{numhypRed} Domain Subdivision}\\
\begin{itemize}
\item {\tt chebfun2} divides the domain to reduce the matrix polynomial problem size \cite[Sec.\ 4]{nakatsukasa_2013}
\item We have yet to implement a domain subdivision method
\end{itemize}
\end{block}
\begin{block}{Current Results}
\begin{center}
\begin{tabular}{ |c|c|c|c|c|  }
 \hline
 & \multicolumn{2}{|c|}{{\tt ChebTools}}&\multicolumn{2}{|c|}{{\tt chebfun}} \\
 \hline
 Functions							& $\lVert\cdot\rVert_2$ Error 	& Time (s)	& $\lVert\cdot\rVert_2$ Error	& Time (s)\\
 \hline
 $F_1(x,y),F_2(x,y)$ & $6.97\times10^{-32}$&0.004268 &$7.7\times10^{-16}$ &0.522 \\
 \hline
 $G_1(x,y),G_2(x,y)$ & $6.21\times10^{-16}$ & 214.059 &$2.92\times10^{-10}$ &0.294 \\
 \hline
 $H_1(x,y),H_2(x,y)$ &$3.24\times10^{-16}$ & 195.903 &$2.77\times10^{-11}$ &0.296 \\
 \hline
\end{tabular}
\end{center}
\begin{align*}
F_1(x,y) &= T_3(x)-13T_1(x) & F_2(x,y) &= T_3(y)-13T_1(y)\\
G_1(x,y) &= \cos(\pi x)(y-2) & G_2(x,y) &= (y-.9)(x-2)\\
H_1(x,y) &= \cos\left(\pi x-\frac{\pi}{10}\right)(y-2) & H_2(x,y) &= (y-.1)(y-.9)(x-2)
\end{align*}
{\color{numhypRed} Discussion}
\begin{itemize}
\item {\tt chebfun} is much faster than {\tt ChebTools}
\item This speed difference can be explained by the lack of a subdivision strategy in our implementation
\item A new subdivision strategy is part of current and future work
\end{itemize}
\end{block}
\end{column}
%% end of second column %%%%%%%%%%%%%%%%%%%%%%%%%%%%%%%%%%%%%%%%%%%%%%%%%%%%%%%%%%%%


%%%% begining of third column%%%%%%%%%%%%%%%%%%%%%%%%%%%%%%%%%%%%%%%%%%%%%%%%%%%%%%%%%%%%
  \begin{column}{.31\linewidth}

%\begin{block}{Current Results}
%\begin{tabular}{ |c|c|c|c|c|  }
% \hline
% & \multicolumn{2}{|c|}{{\tt ChebTools}}&\multicolumn{2}{|c|}{{\tt chebfun}} \\
% \hline
% Functions							& 2 Norm Error 	& Time (s)	& 2 Norm Error	& Time (s)\\
% \hline
% $F_1(x,y),F_2(x,y)$ & $6.97\times10^{-32}$&0.004268 &$7.7\times10^{-16}$ &0.522 \\
% \hline
% $G_1(x,y),G_2(x,y)$ & $6.21\times10^{-16}$ & 214.059 &$2.92\times10^{-10}$ &0.294 \\
% \hline
% $H_1(x,y),H_2(x,y)$ &$3.24\times10^{-16}$ & 195.903 &$2.77\times10^{-11}$ &0.296 \\
% \hline
%\end{tabular}
%\begin{align*}
%F_1(x,y) &= T_3(x)-13T_1(x) & F_2(x,y) &= T_3(y)-13T_1(y)\\
%G_1(x,y) &= \cos(\pi x)(y-2) & G_2(x,y) &= (y-.9)(x-2)\\
%H_1(x,y) &= \cos\left(\pi x-\frac{\pi}{10}\right)(y-2) & H_2(x,y) &= (y-.1)(y-.9)(x-2)
%\end{align*}
%\begin{itemize}
%\item {\tt chebfun} is much faster than {\tt ChebTools}
%\item {\tt chebfun} has a domain subdivision strategy for reducing the size of the generalized eigenvalue problem
%\end{itemize}
%
%\end{block}



\begin{block}{Current Work}
{\color{numhypRed} Our Application Motivates Current Work}
\begin{itemize}
\item Empirical equations of state in thermodynamics have ranges of validity which are not rectangular domains
\item Have to use coordinate transformations \cite{townsend_2013-2} or to work on non-rectangular domains in {\tt chebfun2} 
\end{itemize}
{\color{numhypRed} Our Subdivision Idea}
\begin{itemize}
\item Subdivide where domain boundary is horizontal, vertical, or has a derivative discontinuity
\item Interior subdomains can proceed with rectangular subdivision like in the {\tt chebfun} algorithm
\item Approximate on exterior subdomains using interior Chebyshev points
\end{itemize}

{\color{numhypRed} Example Domain}\\
Interior of the closed curve $(1.5\cos t + .15\sin2t, \sin t + .3\cos t)$ with $t\in[0,2\pi]$

\begin{figure}
\captionsetup[subfigure]{justification=centering,position=bottom}
%\centering
\begin{minipage}{.5\linewidth}
\vspace*{6mm}
\subcaptionbox{Boundary in black and initial subdivisions in blue.}{
\resizebox {\linewidth} {!} {
\begin{tikzpicture}
%\tikzmath{\a = 1.5; \b =.15; \c = 1; \d = .3; }
\def\a{1.5}
\def\b{.15}
\def\c{1}
\def\d{.3}

\def\tvv{rad(asin((-\a+(\a^2+32*(\b^2))^.5)/(8*\b)))}
\def\tv{\tvv+pi}

\def\xv{\a*(cos(deg(\tv))) + \b*(sin(deg(2*\tv)))}
\def\yv{-0.108616270206} %determined numerically

\def\xvv{\a*(cos(deg(\tvv))) + \b*(sin(deg(2*\tvv)))}
\def\yvv{\c*(sin(deg(\tvv))) + \d*(cos(deg(\tvv)))}

\def\th{rad(atan(\c/\d))}
\def\thh{rad(atan(\c/\d))+pi}

\def\xh{\a*(cos(deg(\th))) + \b*(sin(deg(2*\th)))}
\def\yh{\c*(sin(deg(\th))) + \d*(cos(deg(\th)))}


\def\xhh{-0.34845302101} %determined numerically
\def\yhh{\c*(sin(deg(\thh))) + \d*(cos(deg(\thh)))}
\def\ep{.0001}
\begin{axis}[
  xmin=\xv-\ep,xmax=\xvv+\ep,
  ymin=\yhh-\ep,ymax=\yh+\ep,
  x label style={at={(axis description cs:0.5,-0.04)},anchor=north},
  y label style={at={(axis description cs:0.04,.5)},anchor=south},
  yticklabel style={text width = .5cm, anchor=west, align=left,xshift = -.25cm,yshift=0cm},
  yticklabel pos = right,
  xticklabel style={anchor=north},
  xticklabel shift=-8pt
]
  \addplot+[line width=1.5pt,mark = none, black, domain=0:2*pi,samples=100,variable=\t](%
    { \a*(cos(deg(t))) + \b*(sin(deg(2*t))) },%
    {\c*(sin(deg(t))) + \d*(cos(deg(t)))}%
  ); 
  
  
  \addplot+[line width=1.5pt,mark = none, blue, domain=\yh:\yhh,variable=\t]({\xh},{t});
  \addplot+[line width=1.5pt,mark = none, blue, domain=\yh:\yhh,variable=\t]({\xhh},{t});
  \addplot+[line width=1.5pt,mark = none, blue, domain=\xvv:\xv,variable=\t]({t},{\yvv});
  \addplot+[line width=1.5pt,mark = none, blue,, domain=\xvv:\xv,variable=\t]({t},{\yv});
\end{axis}
\end{tikzpicture}}}
\end{minipage}%
\begin{minipage}{.5\linewidth}

\subcaptionbox{Interior 7x7 Chebyshev nodes in red.}{
\resizebox {\linewidth} {!} {
\begin{tikzpicture}


\def\a{1.5}
\def\b{.15}
\def\c{1}
\def\d{.3}

\def\tvv{rad(asin((-\a+(\a^2+32*(\b^2))^.5)/(8*\b)))}
\def\tv{\tvv+pi}

\def\xv{\a*(cos(deg(\tv))) + \b*(sin(deg(2*\tv)))}
\def\yv{-0.108616270206} %determined numerically

\def\xvv{\a*(cos(deg(\tvv))) + \b*(sin(deg(2*\tvv)))}
\def\yvv{\c*(sin(deg(\tvv))) + \d*(cos(deg(\tvv)))}

\def\th{rad(atan(\c/\d))}
\def\thh{rad(atan(\c/\d))+pi}

\def\xh{\a*(cos(deg(\th))) + \b*(sin(deg(2*\th)))}
\def\yh{\c*(sin(deg(\th))) + \d*(cos(deg(\th)))}


\def\xhh{-0.34845302101} %determined numerically
\def\yhh{\c*(sin(deg(\thh))) + \d*(cos(deg(\thh)))}
\def\ymid{.5*(\yh + \yvv)}
\def\ystep{.5*(\yh - \yvv)}

\def\xmid{.5*(\xh + \xvv)}
\def\xstep{.5*(\xvv - \xh)}
\def\ep{.0001}
\begin{axis}[
  xmin=\xh-\ep,xmax=\xvv+\ep,
  ymin=\yvv-\ep,ymax=\yh+\ep,
  x label style={at={(axis description cs:0.5,-.225)},anchor=south},
  y label style={at={(axis description cs:.04,.5)},anchor=south, text height = .01cm},
  ticklabel style={font=\tiny,},
  yticklabel style={anchor=west, align=left},
  yticklabel pos = right,
  yticklabel shift = -10pt,
  xticklabel style={anchor=north},
  xticklabel shift=-8pt
]
  \addplot+[line width=1.5pt,mark = none, black, domain=0:pi/2,samples=100,variable=\t](%
    { \a*(cos(deg(t))) + \b*(sin(deg(2*t))) },%
    {\c*(sin(deg(t))) + \d*(cos(deg(t)))}%
  ); 
   \addplot+[line width=1.5pt,mark = none, blue, domain=\yh:\yhh,variable=\t]({\xh},{t});
  \addplot+[line width=1.5pt,mark = none, blue, domain=\xvv:\xv,variable=\t]({t},{\yvv});
  
  
  \addplot[red, only marks, mark=*] coordinates {
(1.52865090951, 0.480897593475)

(1.46065476435, 0.621680857829)
(1.46065476435, 0.518620355466)
(1.46065476435, 0.480897593475)

(1.27488584106, 0.762464122183)
(1.27488584106, 0.621680857829)
(1.27488584106, 0.518620355466)
(1.27488584106, 0.480897593475)

(1.0211207726, 0.903247386537)
(1.0211207726, 0.762464122183)
(1.0211207726, 0.621680857829)
(1.0211207726, 0.518620355466)
(1.0211207726, 0.480897593475)

(0.767355704145, 1.0063078889)
(0.767355704145, 0.903247386537)
(0.767355704145, 0.762464122183)
(0.767355704145, 0.621680857829)
(0.767355704145, 0.518620355466)
(0.767355704145, 0.480897593475)

(0.581586780849, 1.0063078889)
(0.581586780849, 0.903247386537)
(0.581586780849, 0.762464122183)
(0.581586780849, 0.621680857829)
(0.581586780849, 0.518620355466)
(0.581586780849, 0.480897593475)
(0.513590635689, 1.04403065089)
(0.513590635689, 1.0063078889)
(0.513590635689, 0.903247386537)
(0.513590635689, 0.762464122183)
(0.513590635689, 0.621680857829)
(0.513590635689, 0.518620355466)
(0.513590635689, 0.480897593475)};
\end{axis}
\end{tikzpicture}}}
\end{minipage}
\caption{Example domain on the left and the upper right subdomain on the right}
\end{figure}
\end{block}

\begin{block}{Future Work}
\begin{itemize}
\item Parallelize the {\tt ChebTools} library
\item Adaptive capabilities for approximations
\end{itemize}
Add Features To Construct a Chebyshev expansion from:
\begin{itemize}
\item linear least squares
\item linear boundary value problems
\end{itemize}
\end{block}

\begin{block}{References and Acknowledgements}
\bibliographystyle{plain}
\scriptsize{\bibliography{jmm2018_poster}}
\bigskip
{\normalsize This work was supported by the NIST SURF program.}
\end{block}
\end{column}
  \end{columns}
\end{frame}

\end{document}


%%%%%%%%%%%%%%%%%%%%%%%%%%%%%%%%%%%%%%%%%%%%%%%%%%%%%%%%%%%%%%%%%%%%%%%%%%%%%%%%%%%%%%%%%%%%%%%%%%%%
%%% Local Variables:
%%% mode: latex
%%% TeX-PDF-mode: t
%%% End: 